\section{Introduction}
\label{sec:introduction}

A two-way nondeterministic finite automaton (2NFA) has a finite set of
states and a single read-only head, which may move in either direction
along an input bounded by endmarkers. It accepts when some finite
computation reaches an accepting state. The complementation problem asks
whether every $n$-state 2NFA over a finite alphabet $\Sigma$ has a 2NFA for
the complementary language with at most $p(n)$ states, for one polynomial
$p$ independent of $\Sigma$. We answer this question negatively.

For exact state counts, there is one initial state and a set of accepting
states. Each transition depends only on the current state and scanned
symbol, updates the state, and moves the head left, right, or not at all.
We start the head on the left endmarker and forbid moves beyond either
endmarker. Missing transitions
reject that branch; infinite nonaccepting computations do not accept.
An accepting initial configuration counts as acceptance. All states,
including initial and accepting states, are counted. A deterministic
two-way automaton (2DFA) has at most one successor for each state and
scanned symbol.

For a finite set $H$, let $\mathcal R_H$ denote all binary relations on
$H$. Products are composed in path order:
$(p,r)\in AB$ if and only if $(p,q)\in A$ and $(q,r)\in B$ for some $q$.
Use $\Sigma_H=\mathcal R_H$ as an input alphabet. The product $r(w)$ of a
word $w=R_1\cdots R_k$ is $R_1\cdots R_k$, and the empty product is
$r(\varepsilon)=I_H=\{(p,p):p\in H\}$. Define the relation-product
liveness language
\[
 L_H=\{w\in\Sigma_H^*:r(w)\ne\varnothing\}.
\]
In particular, $L_H$ contains the empty word when $H$ is nonempty.
For nonempty words, this is Sakoda and Sipser's family $B_h$, where
$h=|H|$ \cite[Section~2.1 and Theorem~2.3]{SakodaSipser1978}; the name
\emph{one-way liveness} is used in \cite[Section~2.2]{Kapoutsis2006}.

\begin{theorem}\label{main:theorem}
For every integer $n\ge4$, set $H=\{1,\ldots,n-2\}$ and
$\Sigma_n=\mathcal R_H$. There is an $n$-state 2NFA $A_n$ recognizing
$L_H$ such that every 2NFA recognizing
$\Sigma_n^*\setminus L(A_n)$ has at least
\[
 \frac12\,2^{\lfloor(n-4)/127\rfloor}-1
\]
states. In particular, there is no polynomial state bound for
complementation that is independent of the finite alphabet.
\end{theorem}

The alphabet $\Sigma_n$ consists of all binary relations on a set of
$n-2$ points, so $|\Sigma_n|=2^{(n-2)^2}$. The parameter is the
number of states, not the size of the transition table. Theorem~\ref{main:theorem}
concerns bounds uniform over finite alphabets; it does not assert a lower
bound over one fixed alphabet.

\paragraph{History and significance.}
Sakoda and Sipser's work on nondeterminism and two-way finite automata
\cite{SakodaSipser1978} initiated the central state-succinctness problem of
simulating two-way nondeterministic automata deterministically.
Complementation is a separate question: here the target remains
nondeterministic. Vardi's construction gives an exponential-state
one-way nondeterministic automaton for the complement of a two-way
nondeterministic automaton \cite[Theorem~3.2]{Vardi1989}. It certifies
nonacceptance by locally consistent sets of states. For sweeping
automata, whose head can reverse direction only at the endmarkers,
Kapoutsis proved that the complement of one-way liveness requires
exponentially many states in every sweeping 2NFA
\cite[Theorem~1]{Kapoutsis2006}. Theorem~\ref{main:theorem}
allows unrestricted two-way motion in the complementing automaton.

For deterministic two-way automata, complementation
has a linear state bound independent of the alphabet
\cite{GeffertMereghettiPighizzini2007}.
Guillon, Prigioniero, and Taheri describe polynomial 2NFA
complementation as open \cite[Section~1]{GuillonPrigionieroTaheri2026}
and give polynomial complementation by 1-limited automata
\cite[Theorem~4.1]{GuillonPrigionieroTaheri2026}. These automata may
rewrite a tape cell on its first visit; this resource is absent from
the read-only model considered here.
Theorem~\ref{main:theorem} resolves the alphabet-uniform polynomial
complementation question negatively. The same growing-alphabet family
also requires exponentially many states in every equivalent deterministic
two-way automaton: a smaller deterministic simulator could be
complemented with linear state overhead. Corollary~\ref{auto:determinization}
makes this implication explicit, including the treatment of infinite
nonaccepting computations.

The companion paper \cite[Theorem~1.1]{CompanionDeterminization2026}
proves a deterministic lower bound with a larger exponential rate for the
same language, by an independent matching-diagram argument. We compare
the exact bounds after Corollary~\ref{auto:determinization}. The new
obstruction here applies even when the complementing machine is
nondeterministic.

\paragraph{The argument.}
An automaton with $h+2$ states guesses a path through the
relations and accepts exactly when their product is nonempty.
We prove that recognizing product emptiness requires exponentially many
states.

The proof separates a representation of computations from a finite
algebraic lower bound. A segment of a computation is represented by four
relations recording the possible entries and exits through its two
boundaries. These \emph{path diagrams} compose by joining boundaries.
Their monoid is the monoid of partitioned binary relations introduced
by Martin and Mazorchuk \cite[Sections~2.1--2.3]{MartinMazorchuk2013},
written with separate labels for the two directions of travel.
Adding edges only increases the available accepting paths. On the other
hand, singleton relation contexts test the absence of any prescribed
pair. Consequently a machine for product emptiness induces a
surjective multiplicative map from its path diagrams to the relation
monoid that reverses inclusion.

To bound such maps, we attach recurrent classes of boundary labels to an
idempotent diagram $e$, meaning $e^2=e$. These classes group labels joined
in both directions by through paths together with returns. Their precise
definition is given in Section~\ref{sec:diagrams}.
The diagrams $z$ satisfying $ez=ze=z$ form a
submonoid with identity $e$, called its \emph{corner}. Such a $z$ may
destroy the loops on some of those classes; we record the destroyed
classes as a missing set. Two structural facts
control this loss. First, one common missing-set bound transports through
any fixed diagram context with only a factor of two. Second, along a
chain of nested idempotents the sum of successive losses is bounded by
twice a single initial missing-set budget. The second fact permits the
corner identity to change during the argument.

For an integer parameter $t$, the amplification step converts $2t$
conjugates of one relation with an added off-diagonal pair
into $t^2$ additions sharing a common budget. Each addition, after
restriction to a smaller relation corner, incurs an inductively bounded
loss. Taking $t=64$ gives the exponential recurrence. This strategy
adapts the conjugate-amplification pattern of the matching-diagram
rank-loss argument in
\cite[Section~3]{CompanionDeterminization2026}; that argument concerns a deterministic
state lower bound. Here all diagram paths may be nondeterministic, and
the transport and nested-class bounds are proved anew. They are stated
independently of automata and may be useful for other order-reversing
images of finite path monoids.

Section~\ref{sec:diagrams} states the algebraic bound and derives the
order-reversing map before developing recurrent classes and transport.
Section~\ref{nest:section} proves the chain bound, and
Section~\ref{amp:section} proves the algebraic lower bound.
Section~\ref{auto:section} proves the promised automata representation and
completes the proof of Theorem~\ref{main:theorem}.

The bounds impose no restriction on input length or on the running time
of a successful computation. They compare finite-state machines without
a work tape, and give no separation between uniform logarithmic-space
complexity classes.
